% !TeX root = ../manuscript.tex

% 中英文摘要和关键字

\begin{abstract}
  摘要是学位论文的内容不加注释和评论的简要陈述。
  摘要应具有独立性和自含性，即不阅读论文的全文，就能获得必要的信息。
  摘要一般应说明研究工作的目的、方法、结果和结论等，
  应突出论文的创新点。
  
  博士学位论文的摘要一般2000字左右，硕士学位论文的摘要一般1000字左右。
  摘要中应尽量避免采用图、表、化学结构式、非公知公用的符号和术语。
  英文摘要的内容应与中文摘要相对应。

  关键词是为了便于文献索引和检索工作，
  从学位论文中选取出来用以表示全文主题内容信息的单词或术语。
  每篇论文选取3－8个关键词，每个关键词之间用逗号间隔，另起一行，排在摘要内容下方。

  % 关键词用“英文逗号”分隔，输出时会自动处理为正确的分隔符
  \bnusetup{
    keywords = {关键词 1, 关键词 2, 关键词 3, 关键词 4, 关键词 5},
  }
\end{abstract}

\begin{abstract*}
An abstract is a brief statement of the content of a thesis without annotations or comments. The abstract should be independent and self-contained, meaning that necessary information can be obtained without reading the full text of the paper. The abstract should generally explain the purpose, methods, results, and conclusions of the research work, and highlight the innovative points of the paper.

The abstract for a doctoral thesis is generally around 2000 words, while the abstract for a master's thesis is generally around 1000 words. The use of diagrams, tables, chemical structural formulas, symbols and terminology that are not commonly used should be avoided as much as possible in the abstract. The content of the English abstract should correspond to the Chinese abstract.

Keyword is a word or term selected from a thesis to represent the topic information of the entire text for the convenience of literature indexing and retrieval work. Select 3-8 keywords for each paper, separated by commas, and start a new line below the abstract content.

  % Use comma as separator when inputting
  % 英文关键词首字母请大写
  \bnusetup{
    keywords* = {Keyword 1, Keyword 2, Keyword 3, Keyword 4, Keyword 5},
  }
\end{abstract*}
